\documentclass[11pt,oneside]{article}

\usepackage[a4paper,margin=27mm,headheight=15pt]{geometry}
\usepackage[T1]{fontenc}
\usepackage[utf8]{inputenc}
\IfFileExists{libertinus.sty}{\usepackage{libertinus}}{\usepackage{lmodern}}
\usepackage{microtype}
\usepackage{amsmath,amssymb,amsthm,mathtools,mathrsfs}
\usepackage{bm}
\usepackage{booktabs,longtable,array,tabularx}
\usepackage{graphicx}
\usepackage{pgfplots}
\usepackage{xcolor}
\usepackage{enumitem}
\usepackage{fancyhdr}
\usepackage{titlesec}
\usepackage{listings}
\usepackage{xurl}
\usepackage{hyperref}
\usepackage[nameinlink,noabbrev]{cleveref}

\definecolor{linkblue}{RGB}{25,75,135}
\definecolor{shadegray}{RGB}{246,247,249}
\definecolor{rulegray}{RGB}{195,201,208}
\pgfplotsset{compat=1.18}
\hypersetup{
  colorlinks=true,
  linkcolor=linkblue,
  citecolor=linkblue,
  urlcolor=linkblue,
  pdftitle={Small-argument asymptotics research frontiers for the Fabius function},
  pdfsubject={Coefficient formulas, numerical evidence, and asymptotic claims awaiting formalization},
  pdfkeywords={Fabius function, asymptotic expansion, Lambert W, saddle point, Stirling numbers, Gamma-zeta Fourier coefficients}
}

\pagestyle{fancy}
\fancyhf{}
\fancyhead[L]{\small Small-argument asymptotics research frontiers}
\fancyhead[R]{\small\nouppercase{\leftmark}}
\fancyfoot[C]{\thepage}
\renewcommand{\headrulewidth}{0.4pt}

\titleformat{\section}{\Large\bfseries}{\thesection}{0.7em}{}
\titleformat{\subsection}{\large\bfseries}{\thesubsection}{0.7em}{}
\titleformat{\subsubsection}{\normalsize\bfseries}{\thesubsubsection}{0.7em}{}
\setcounter{secnumdepth}{3}
\setcounter{tocdepth}{2}
\setlist[itemize]{leftmargin=2em,itemsep=0.25em,topsep=0.4em}
\setlist[enumerate]{leftmargin=2.3em,itemsep=0.25em,topsep=0.4em}

\newtheorem{theorem}{Theorem}[section]
\newtheorem{proposition}[theorem]{Proposition}
\newtheorem{lemma}[theorem]{Lemma}
\newtheorem{corollary}[theorem]{Corollary}
\newtheorem{conjecture}[theorem]{Conjecture}
\theoremstyle{definition}
\newtheorem{definition}[theorem]{Definition}
\newtheorem{algorithm}[theorem]{Algorithm}
\newtheorem{example}[theorem]{Example}
\theoremstyle{remark}
\newtheorem{remark}[theorem]{Remark}
\newtheorem{warning}[theorem]{Warning}

\newcommand{\R}{\mathbb R}
\newcommand{\C}{\mathbb C}
\newcommand{\Q}{\mathbb Q}
\newcommand{\N}{\mathbb N}
\newcommand{\Z}{\mathbb Z}
\newcommand{\Up}{\operatorname{up}}
\newcommand{\sinc}{\operatorname{sinc}}
\newcommand{\wt}{\operatorname{wt}_2}
\newcommand{\e}{\mathrm e}
\newcommand{\ii}{\mathrm i}
\newcommand{\dd}{\,\mathrm d}
\newcommand{\Li}{\operatorname{Li}}
\newcommand{\Var}{\operatorname{Var}}
\newcommand{\E}{\mathbb E}
\newcommand{\Prob}{\mathbb P}
\newcommand{\bigO}{\mathcal O}
\newcommand{\smallo}{o}
\newcommand{\supp}{\operatorname{supp}}
\newcommand{\sgn}{\operatorname{sgn}}
\newcommand{\lcm}{\operatorname{lcm}}
\newcommand{\vtwo}{v_2}
\newcommand{\qbinom}[3]{\genfrac{[}{]}{0pt}{}{#1}{#2}_{#3}}
\newcommand{\repo}{\nolinkurl{Analysis/FabiusFunction}}
\newcommand{\Psihat}{\widehat\Psi}

\newenvironment{proofidea}{\par\smallskip\noindent\textit{Proof architecture.}\ }{\hfill$\square$\par\smallskip}
\newenvironment{boxedremark}{%
  \par\smallskip\noindent\begin{center}\begin{minipage}{0.94\textwidth}%
  \color{black}\hrule\vspace{0.6em}\small
}{%
  \vspace{0.6em}\hrule\end{minipage}\end{center}\smallskip
}

\lstdefinestyle{pseudo}{
  basicstyle=\ttfamily\small,
  backgroundcolor=\color{shadegray},
  frame=single,
  rulecolor=\color{rulegray},
  columns=fullflexible,
  keepspaces=true,
  showstringspaces=false,
  xleftmargin=0.7em,
  xrightmargin=0.7em,
  aboveskip=0.8em,
  belowskip=0.8em
}

\title{\bfseries Small-Argument Asymptotics Research Frontiers\\[0.35em]
\large Coefficient formulas and numerical evidence awaiting formalization}
\author{Research companion to the formal development\\
\small Vladimir Reshetnikov's \texttt{ProveIt} repository}
\date{Repository snapshot: 24 August 2026}

\begin{document}
\maketitle

\begin{abstract}
As $x\to0^+$ the Fabius function $F$ is flat to infinite order, so it has no
Taylor expansion at $0$; the correct description is an exponentially small
asymptotic expansion in the lower-Lambert phase $\lambda(x)$, whose
coefficients $A_j$ are one-periodic functions of $\lambda$ rather than
constants.  Until now those coefficients existed only as the output of a chain
of recursions --- a differential recursion for the saddle jets, a formal
exponential recursion, a Gaussian contraction, and a formal logarithm --- and
only $A_1$ had ever been written down.  This note closes the chain.  We solve
the differential recursion in closed form using shifted signed Stirling numbers
of the first kind and harmonic numbers, collapse the exponent coefficients to a single
power of $\ii$ times a jet term plus a universal jet-free tail, and assemble a
completely explicit, non-recursive formula for every $A_j$ as a finite sum over
compositions, subsets and Gaussian moments.  As applications we compute $A_2$
and $A_3$ explicitly for the first time, identify the highest-derivative term
of $A_j$ as $(-1)^j\Psi^{(2j)}/(2^jj!\,L^{2j})$, deduce the
Gamma--zeta Fourier form and the amplitude law of every coefficient, and verify
the whole chain numerically against exact rational values of $F(2^{-n})$.
\end{abstract}

\begin{boxedremark}
\textbf{Formalization boundary.}
This document preserves the former small-argument draft as a research-frontier notebook.
The sharp main term, recursive all-orders expansion, closed product-polynomial saddle jets,
and the first two corrections have exact Lean counterparts, but several stronger claims in
the pages below do not.  In particular, the ordered-composition formulas, identification
with standard shifted Stirling numbers, general logarithmic highest-derivative law,
expanded higher coefficients, Fourier-amplitude estimates, uniqueness and phase-locking
claims, and all numerical tables remain to be formalized.  Those passages are research
targets, not established results of the project.
\end{boxedremark}

\tableofcontents

\section{Introduction}

Throughout, $F$ is the canonical bounded Fabius function on $[0,1]$: the
cumulative distribution function of
\[
  X=\sum_{k\ge1}2^{-k}U_k,\qquad U_1,U_2,\dots\ \text{independent uniform on }[0,1].
\]
It is smooth, vanishes on $(-\infty,0]$, equals $1$ on $[1,\infty)$, satisfies
$F(x)+F(1-x)=1$, and obeys $F'(x)=2F(2x)$ on $[0,\tfrac12]$.  Every derivative
vanishes at the origin, so $F$ is flat there and no Taylor expansion exists.

This note develops the behaviour of $F(x)$ as $x\to0^+$ in detail, and in
particular gives a general formula for every coefficient of its all-orders
expansion.  It is a companion to the formal development in \repo{} and to the
article in \nolinkurl{docs/Fabius_Function_and_Rvachev_Up/}, and it is
self-contained: everything needed to compute any coefficient by hand is stated
here.

\subsection{What was already known, and what is new}

The pre-existing development already contained the corrected sharp main term
$\mathcal M(x)$ with its genuine nonconstant periodic correction, a proof that
deleting that correction makes the standard printed formula false, and an
all-orders expansion
\[
  \log F(x)=\mathcal M(x)+\sum_{j\ge1}\frac{A_j(\lambda)}{\lambda^j}
\]
in which the coefficient functions $A_j$ were \emph{defined only by a
recursion}.  Only $A_1$ was ever displayed.

What this note adds:

\begin{enumerate}
\item A closed form for the saddle jets (\cref{sec:jets}).  The differential
  recursion is solved: the $n$-th jet is a harmonic-number constant plus a
  Stirling-weighted combination of the first $n+1$ derivatives of the periodic
  correction.  This removes the only non-algebraic recursion from the chain.
\item A closed form for the exponent coefficients (\cref{sec:exponent}), with
  the factorials and the spurious extra power of $\ii$ cancelled.
\item A fully explicit, non-recursive formula for every coefficient
  (\cref{sec:general}): $A_j$ is a finite triple sum over compositions, subsets
  and Gaussian moments.
\item The explicit coefficients $A_2$ and $A_3$ (\cref{sec:explicit}), which did
  not exist anywhere before, in two equivalent forms.
\item The leading-derivative law (\cref{sec:leading}): the top term of $A_j$ is
  $(-1)^j\Psi^{(2j)}(u)/(2^jj!\,L^{2j})$, which explains why the oscillation of
  the coefficients grows with their order.
\item The Gamma--zeta Fourier form of every coefficient and the resulting
  amplitude law (\cref{sec:fourier}).
\item An independent end-to-end numerical verification (\cref{sec:numerics})
  against exact rational values of $F(2^{-n})$.
\end{enumerate}

\subsection{Notation}

\begin{center}
\begin{tabular}{@{}ll@{}}
\toprule
symbol & meaning\\
\midrule
$L$ & $\log2=0.69314718055994530942\dots$\\
$\gamma$ & Euler--Mascheroni constant\\
$\gamma_1$ & first Stieltjes constant, $\zeta(1+s)=1/s+\gamma-\gamma_1s+\bigO(s^2)$\\
$H_n$ & $n$-th harmonic number, $H_0=0$\\
$\Psi$ & the centered one-periodic correction of \cref{sec:main}\\
$\Psi^{(m)}$ & $m$-th derivative of $\Psi$\\
$\lambda=\lambda(x)$ & the lower-Lambert phase of \cref{sec:lambert}\\
$\mathcal M(x)$ & the sharp main term of \cref{sec:main}\\
$A_j$ & the $j$-th expansion coefficient, a one-periodic function\\
$a_j$ & the $j$-th \emph{mass} coefficient, before taking the logarithm\\
$d_n$ & the $n$-th bounded exponent jet of \cref{sec:jets}\\
$p_n$ & the $n$-th periodic jet of \cref{sec:jets}\\
$s_n$ & $(-1)^{n+1}n!$, the jet slope\\
$c(n,m)$ & $[z^m]\prod_{k=1}^n(z-k)=s(n+1,m+1)$, a shifted signed Stirling number\\
$(N-1)!!$ & $1\cdot3\cdot5\cdots(N-1)$, the $N$-th standard Gaussian moment for even $N$\\
\bottomrule
\end{tabular}
\end{center}

$Z$ denotes a standard Gaussian variable, so $\E[Z^N]=(N-1)!!$ for even $N$ and
$0$ for odd $N$.

\section{The lower-Lambert phase}\label{sec:lambert}

The saddle point of the Laplace inversion sits at radius $r=2^\lambda$, where
$\lambda$ solves
\begin{equation}\label{eq:lambda}
  \lambda\,2^{-\lambda}=x,
  \qquad\text{equivalently}\qquad
  L\lambda=-\log x+\log\lambda .
\end{equation}
This is the lower branch of the Lambert $W$ function,
$\lambda(x)=-W_{-1}(-Lx)/L$, defined and strictly decreasing for
$0<Lx<\e^{-1}$, with $\lambda(x)\to+\infty$ as $x\to0^+$.

The expansion below is an expansion in $1/\lambda$, \emph{not} in $1/(-\log x)$.
The two agree to leading order, since
$L\lambda=-\log x+\log(-\log x)-\log L+\smash{o(1)}$, but they differ at every
order after that; \cref{sec:notclaimed} explains why substituting an expansion
of $\lambda$ into the coefficients is not a legitimate step.

For dyadic arguments $x=2^{-n}$, equation \cref{eq:lambda} reads
$L\lambda=nL+\log\lambda$, whose solution is conveniently obtained by the
iteration $\lambda\mapsto n+\log(\lambda)/L$.

\section{The periodic correction and the sharp main term}\label{sec:main}

Let
\[
  G(-s)=\prod_{j\ge1}\frac{1-\e^{-s/2^j}}{s/2^j},
  \qquad \Lambda(s)=\log G(-s),
\]
be the logarithm of the two-sided Laplace transform of the Fabius density.
Removing the quadratic drift and the forward tail
$T(s)=\sum_{m\ge0}\log(1-\e^{-s2^m})$ produces the \emph{exactly} one-periodic
function
\[
  P(t)=\Lambda(2^t)+\frac L2\bigl(t^2-t\bigr)+T(2^t),
  \qquad P(t+1)=P(t),
\]
so that
\[
  \Lambda(s)=-\frac{(\log s)^2}{2L}+\frac{\log s}{2}+P(\log_2s)+E(s),
  \qquad E=-T,\qquad
  0\le E(s)\le\frac{\e^{-s}}{(1-\e^{-s})^2}.
\]
Write $\Psi(t)=P(t)-\bar P$ with $\bar P=\int_0^1P$.  Then $\Psi$ is
one-periodic, $C^\infty$, of mean zero, and --- this is the decisive point ---
\emph{not constant}.  Its Fourier coefficients are
\begin{equation}\label{eq:psihat}
  \Psihat(0)=0,\qquad
  \Psihat(k)=-\frac{\Gamma(-\chi_k)\,\zeta(1-\chi_k)}{L},
  \qquad \chi_k=\frac{2\pi\ii k}{L},\quad k\ne0,
\end{equation}
which are nonzero for every $k\ne0$ because $\Gamma$ has no zeros and $\zeta$
has none on the line $\operatorname{Re}s=1$.  Numerically
\[
\begin{aligned}
  \Psihat(1)&=(\;7.55864754098-7.47010355898\,\ii\;)\times10^{-7},
    &|\Psihat(1)|&=1.06271162519\times10^{-6},\\
  \Psihat(2)&=(\;3.24812347723+5.85175898712\,\ii\;)\times10^{-13},
    &|\Psihat(2)|&=6.69278636792\times10^{-13},\\
  \Psihat(3)&=(-2.81236523199-2.58674397916\,\ii\;)\times10^{-19},
    &|\Psihat(3)|&=3.82107872359\times10^{-19},
\end{aligned}
\]
so $\Psi$ itself oscillates by only about $2\times10^{-6}$.  \Cref{sec:traps}
records the numerical trap that makes these values easy to get wrong.

With $A_0^{\mathrm{c}}=(6\gamma^2+12\gamma_1-\pi^2)/12$ and
\[
  C_F=\frac{A_0^{\mathrm{c}}}{L}-\frac{7L}{12}-\frac{\log\pi}{2}
     =-2.027983898638859423971\dots,
\]
the sharp main term is
\begin{equation}\label{eq:main}
  \mathcal M(x)=-\frac L2\lambda^2+\Bigl(1+\frac L2\Bigr)\lambda
    -\frac12\log\lambda+C_F+\Psi(\lambda),
\end{equation}
and $\log F(x)=\mathcal M(x)+\bigO(1/\lambda)$.  Deleting $\Psi(\lambda)$
invalidates that error term: the ratio $F(x)/\exp(\mathcal M(x)-\Psi(\lambda))$
has at least two distinct subsequential limits.

\section{The saddle jets and their closed form}\label{sec:jets}

Let $q=\log G(-\cdot\,)$ evaluated on the saddle ray $r=2^t$.  Up to the
exponentially small forward tail, the $(n+1)$-st scaled derivative has the form
\[
  r^{\,n+1}q^{(n+1)}(r)=s_nt+p_n(t),\qquad s_n=(-1)^{n+1}n!,
\]
where the periodic jets $p_n$ are determined by the differential recursion
\begin{equation}\label{eq:jetrec}
  p_0(t)=\frac12+\frac{\Psi'(t)}{L},
  \qquad
  p_{n+1}(t)=\frac{p_n'(t)}{L}-(n+1)p_n(t)+\frac{s_n}{L}.
\end{equation}
The quantity that actually enters the saddle exponent is the \emph{bounded
exponent jet} $d_n=p_n+s_n$.

\subsection{Solving the recursion}

Write $D$ for $\dd/\dd t$.  The recursion \cref{eq:jetrec} reads
$p_{n+1}=\bigl(D/L-(n+1)\bigr)p_n+s_n/L$, and the operators $D/L-j$ commute,
since
\[
  \Bigl(\frac DL-j\Bigr)\Bigl(\frac DL-k\Bigr)
   =\frac{D^2}{L^2}-(j+k)\frac DL+jk
\]
is symmetric in $j$ and $k$.  Hence the homogeneous part of the solution is the
falling-factorial operator $\prod_{k=1}^n(D/L-k)$ applied to $p_0$.  Each
inhomogeneous constant $s_j/L$ is carried forward by
$\prod_{k=j+2}^n(D/L-k)$, and a constant is an eigenvector of $D/L-k$ with
eigenvalue $-k$, so its contribution is
\[
  \frac{s_j}{L}\,(-1)^{\,n-j-1}\frac{n!}{(j+1)!}
   =\frac{(-1)^n\,n!}{(j+1)L}.
\]
Summing over $j=0,\dots,n-1$ produces the harmonic number.  Finally
$\prod_{k=1}^n(z-k)=\sum_{m=0}^nc(n,m)z^m$, so applying the operator polynomial
to $p_0=\tfrac12+\Psi'/L$ gives the result.

\subsection{The closed forms}

\begin{theorem}[Closed-form jets]\label{thm:jets}
For every $n\ge0$ and every $t$,
\[
\begin{aligned}
  p_n(t)&=(-1)^nn!\Bigl(\frac12+\frac{H_n}{L}\Bigr)
    +\sum_{m=0}^{n}c(n,m)\,\frac{\Psi^{(m+1)}(t)}{L^{m+1}},\\
  d_n(t)&=(-1)^nn!\Bigl(\frac{H_n}{L}-\frac12\Bigr)
    +\sum_{m=0}^{n}c(n,m)\,\frac{\Psi^{(m+1)}(t)}{L^{m+1}},
\end{aligned}
\]
where $c(n,m)=[z^m]\prod_{k=1}^n(z-k)$, so that $c(n,n)=1$ and $c(n,m)=0$ for
$m>n$.
\end{theorem}

The coefficient array obeys $c(0,0)=1$, $c(0,m+1)=0$,
$c(n+1,0)=-(n+1)c(n,0)$ and $c(n+1,m+1)=c(n,m)-(n+1)c(n,m+1)$.  The first rows
are:

\begin{center}
\begin{tabular}{@{}ccl@{}}
\toprule
$n$ & $c(n,0),c(n,1),\dots$ & constant part of $d_n$\\
\midrule
$0$ & $1$ & $-1/2$\\
$1$ & $-1,\,1$ & $1/2-1/L$\\
$2$ & $2,\,-3,\,1$ & $-1+3/L$\\
$3$ & $-6,\,11,\,-6,\,1$ & $3-11/L$\\
$4$ & $24,\,-50,\,35,\,-10,\,1$ & $-12+50/L$\\
$5$ & $-120,\,274,\,-225,\,85,\,-15,\,1$ & $60-274/L$\\
\bottomrule
\end{tabular}
\end{center}

\begin{warning}[The index shift]
The weights $c(n,m)$ are shifted signed Stirling numbers of the first kind,
not the unshifted ones.  With the standard convention
$x(x-1)\cdots(x-n+1)=\sum_k s(n,k)x^k$ one has
\[
  z\prod_{k=1}^{n}(z-k)=\prod_{k=0}^{n}(z-k)=\sum_{m}s(n+1,m)z^m,
  \qquad\text{hence}\qquad c(n,m)=s(n+1,m+1).
\]
The distinction is not cosmetic: $s(n,0)=0$ for every $n\ge1$, whereas
$c(n,0)=\prod_{k=1}^n(-k)=(-1)^nn!$, which is the entire constant column of
the table above and carries the harmonic-number part of the jet.  Reading
$c(n,m)$ as $s(n,m)$ would delete it.  Concretely $c(2,\cdot)=(2,-3,1)$ while
$s(2,\cdot)=(0,-1,1)$, and $c(3,\cdot)=(-6,11,-6,1)$ while
$s(3,\cdot)=(0,2,-3,1)$.
\end{warning}

Explicitly, for the first four jets,
\[
\begin{aligned}
  d_0&=\frac{\Psi'}{L}-\frac12,\\
  d_1&=\frac12-\frac1L-\frac{\Psi'}{L}+\frac{\Psi''}{L^2},\\
  d_2&=-1+\frac3L+\frac{2\Psi'}{L}-\frac{3\Psi''}{L^2}+\frac{\Psi'''}{L^3},\\
  d_3&=3-\frac{11}{L}-\frac{6\Psi'}{L}+\frac{11\Psi''}{L^2}
        -\frac{6\Psi'''}{L^3}+\frac{\Psi''''}{L^4}.
\end{aligned}
\]
Since $\Psi$ and all of its derivatives are one-periodic, bounded and of mean
zero, every jet is a bounded one-periodic function whose mean is exactly its
constant part.  This is the source of every ``constant part'' appearing below.

\begin{remark}
\Cref{thm:jets} is formalized as
\lstinline[style=pseudo]{Fabius.negativeLaplacePeriodicJet_eq_closedForm} and
\lstinline[style=pseudo]{Fabius.negativeLaplaceBoundedExponentJet_eq_closedForm}
in \nolinkurl{FabiusSaddleJetClosedForm.lean}.
\end{remark}

\section{The exponent coefficients in closed form}\label{sec:exponent}

After extracting the Gaussian factor and setting $\varepsilon=\lambda^{-1/2}$,
the coefficient of $\varepsilon^m$ in the central saddle exponent is, as
originally recorded,
\[
  E_m(t,v)=\ii^{\,m}\,d_{m-1}(t)\frac{v^m}{m!}
    +\ii^{\,m+2}s_{m+1}\frac{v^{m+2}}{(m+2)!}.
\]
Both the extra power of $\ii$ and both factorials in the second term cancel:
$\ii^{\,m+2}=-\ii^{\,m}$, $s_{m+1}=(-1)^m(m+1)!$ and
$(m+1)!/(m+2)!=1/(m+2)$.

\begin{theorem}[Closed-form exponent]\label{thm:exponent}
$E_0=0$, and for every $m\ge1$,
\[
  E_m(t,v)=\ii^{\,m}\left(d_{m-1}(t)\frac{v^m}{m!}
    +(-1)^{m+1}\frac{v^{m+2}}{m+2}\right).
\]
\end{theorem}

So the second summand is a \emph{universal} tail: it does not see the periodic
correction at all.  The first four are
\[
  E_1=\frac{\ii v(3d_0+v^2)}{3},\quad
  E_2=\frac{v^2(-2d_1+v^2)}{4},\quad
  E_3=\ii v^3\Bigl(-\frac{d_2}{6}-\frac{v^2}{5}\Bigr),\quad
  E_4=\frac{v^4(d_3-4v^2)}{24}.
\]
In particular $E_m(t,-v)=(-1)^mE_m(t,v)$, which is the algebraic reason that
all odd half-powers of $\lambda$ disappear after Gaussian integration, so that
the expansion runs over integer powers $\lambda^{-j}$ only.

\begin{remark}
\Cref{thm:exponent} is formalized as
\lstinline[style=pseudo]{Fabius.negativeLaplaceExponentCoefficient_succ_eq}
and
\lstinline[style=pseudo]{Fabius.negativeLaplaceExponentPolynomial_succ_eq}
in \nolinkurl{FabiusSaddleExponentClosedForm.lean}.
\end{remark}

\section{The general coefficient formula}\label{sec:general}

Write $\mathcal C(N,r)$ for the set of compositions of $N$ into $r$ positive
parts, that is, ordered tuples $(m_1,\dots,m_r)$ of positive integers with
$m_1+\dots+m_r=N$.

\subsection{Step 1: the mass coefficients}

The normalized saddle mass has the expansion
$1+\sum_{j\ge1}a_j(u)\lambda^{-j}$, where $a_j$ is the Gaussian contraction of
the coefficient of $\varepsilon^{2j}$ in
$\exp\bigl(\sum_{m\ge1}E_m\varepsilon^m\bigr)$.  Expanding the exponential over
compositions, and then expanding the product over the two monomials of each
$E_{m_i}$, gives the following.

\begin{theorem}[Explicit mass coefficients]\label{thm:mass}
For every $j\ge1$,
\[
  a_j(u)=(-1)^j\sum_{r=1}^{2j}\frac1{r!}
    \sum_{(m_1,\dots,m_r)\in\mathcal C(2j,r)}\ \sum_{S\subseteq\{1,\dots,r\}}
      \bigl(2j+2|S|-1\bigr)!!
      \prod_{i\in S}\frac{(-1)^{m_i+1}}{m_i+2}
      \prod_{i\notin S}\frac{d_{m_i-1}(u)}{m_i!}.
\]
\end{theorem}

Here $S$ records, for each factor of the product, which of the two monomials of
$E_{m_i}$ was chosen: $i\in S$ picks the universal tail
$v^{m_i+2}/(m_i+2)$, and $i\notin S$ picks the jet term
$d_{m_i-1}v^{m_i}/m_i!$.  The total degree is then $2j+2|S|$, always even, and
its Gaussian moment is the displayed double factorial.  The prefactor $(-1)^j$
is $\ii^{\,2j}$.

\subsection{Step 2: taking the logarithm}

\begin{theorem}[Explicit logarithmic coefficients]\label{thm:log}
With $a_0=1$, for every $j\ge1$,
\[
  A_j(u)=\sum_{r=1}^{j}\frac{(-1)^{r-1}}{r}
    \sum_{(q_1,\dots,q_r)\in\mathcal C(j,r)}a_{q_1}(u)\cdots a_{q_r}(u).
\]
\end{theorem}

\subsection{Step 3: expanding the jets}

Substituting \cref{thm:jets} into \cref{thm:mass,thm:log} expresses $A_j(u)$ as
a finite expression with rational coefficients in $1/L$ and
$\Psi'(u),\dots,\Psi^{(2j)}(u)$, with no recursion anywhere.  Both closed forms
were verified symbolically against the original recursions for $j\le3$, with
identically vanishing residuals.

For reference, the equivalent recursive presentations, which are what the Lean
development uses internally, are
\[
\begin{aligned}
  a_j&=\text{Gaussian contraction of }\operatorname{expCoeff}(E)(2j),\\
  \operatorname{expCoeff}(E)(0)&=1,\qquad
  (n+1)\operatorname{expCoeff}(E)(n+1)
    =\sum_{j\le n}(j+1)E_{j+1}\operatorname{expCoeff}(E)(n-j),\\
  A_0&=0,\qquad A_n=a_n-\frac1n\sum_{k=1}^{n-1}kA_ka_{n-k}.
\end{aligned}
\]

\subsection{Structural consequences}

\begin{itemize}
\item $A_j$ involves $\Psi',\dots,\Psi^{(2j)}$ and no higher derivative,
  because the highest jet reached at order $2j$ is $d_{2j-1}$, which reaches
  $\Psi^{(2j)}$.
\item Every $A_j$ is $C^\infty$, one-periodic and bounded, being a polynomial
  in bounded one-periodic smooth functions.
\item The mean of $A_j$ is its ``$\Psi$-free part'' plus the mean of the
  terms of degree two and higher in the derivatives of $\Psi$.  Those means
  do not vanish, and they are computable in closed form, because
  $\langle\Psi^{(a)}\Psi^{(b)}\rangle$ depends only on $a+b$ up to sign:
  integration by parts over a period gives
  $\langle\Psi'\Psi''\rangle=0$ and
  $\langle\Psi'\Psi'''\rangle=-\langle(\Psi'')^2\rangle$, and generally
  $\langle P_aP_b\rangle=2\sum_{k\ge1}(2\pi\ii k)^a(-2\pi\ii k)^b|\Psihat(k)|^2$,
  which carries the amplification $(2\pi k)^{a+b}$ of \cref{sec:fourier}.
  Applying this to \cref{eq:A2psi} and to $A_1$ leaves
  \begin{align*}
    \langle A_1\rangle&=\frac1{24}+\frac1{2L}
      -\frac{\langle(\Psi')^2\rangle}{2L^2},\\
    \langle A_2\rangle&=\frac1{4L^2}
      -\frac{\langle(\Psi')^2\rangle}{2L^3}
      -\frac{\langle(\Psi'')^2\rangle}{4L^4}
      -\frac{\langle(\Psi')^3\rangle}{6L^3},
  \end{align*}
  every other quadratic and cubic mean cancelling.  Numerically
  $\langle(\Psi')^2\rangle=8.917\times10^{-11}$,
  $\langle(\Psi'')^2\rangle=3.520\times10^{-9}$ and
  $\langle(\Psi')^3\rangle=1.115\times10^{-21}$, so
  $\langle A_1\rangle=0.7630141870184$ and
  $\langle A_2\rangle=0.5203422413049$, against $\Psi$-free parts
  $0.7630141871111$ and $0.5203422452514$.  The gaps are
  $-9.3\times10^{-11}$ and $-3.9\times10^{-9}$, and at $j=3$ the same
  computation gives $-1.14\times10^{-7}$: about a factor $40$ per order,
  for the same reason the oscillation grows.  What justifies calling the
  $\Psi$-free part the mean is therefore not an absolute bound but a ratio:
  at each order the gap sits four to five decades below the amplitude of
  $A_j$ itself.
\end{itemize}

\section{The explicit coefficients}\label{sec:explicit}

It is convenient to state the coefficients first in the jets $d_n$ and then in
the derivatives of $\Psi$.  The jet form is dramatically shorter.

\subsection{Mass coefficients in the jets}

\[
\begin{aligned}
  a_1&=-\frac{6d_0^2+12d_0+6d_1+1}{12},\\[0.3em]
  a_2&=\frac{1}{288}\bigl(36d_0^4+240d_0^3+216d_0^2d_1+300d_0^2+720d_0d_1\\
     &\qquad\qquad+144d_0d_2+24d_0+108d_1^2+300d_1+240d_2+36d_3+1\bigr).
\end{aligned}
\]

\subsection{Logarithmic coefficients in the jets}

\[
  A_1=-\frac{6d_0^2+12d_0+6d_1+1}{12}\;(=a_1),
\]
\begin{equation}\label{eq:A2jets}
  A_2=\frac{1}{24}\bigl(8d_0^3+12d_0^2d_1+12d_0^2+48d_0d_1+12d_0d_2
      +6d_1^2+24d_1+20d_2+3d_3\bigr),
\end{equation}
\begin{equation}\label{eq:A3jets}
\begin{aligned}
  A_3=-\frac{1}{720}\bigl(&180d_0^4+720d_0^3d_1+120d_0^3d_2+240d_0^3
    +360d_0^2d_1^2+2520d_0^2d_1\\
   &+1440d_0^2d_2+180d_0^2d_3+2160d_0d_1^2+720d_0d_1d_2+1440d_0d_1\\
   &+2520d_0d_2+960d_0d_3+90d_0d_4+120d_1^3+1980d_1^2+1800d_1d_2\\
   &+180d_1d_3+150d_2^2+720d_2+780d_3+210d_4+15d_5-2\bigr).
\end{aligned}
\end{equation}

One order further, in the same coordinates,
\begin{equation}\label{eq:A4jets}
\begin{aligned}
  A_4=\frac{1}{5760}\bigl(&1152d_0^5+8640d_0^4d_1+2880d_0^4d_2+240d_0^4d_3+1440d_0^4\\
   &+11520d_0^3d_1^2+2880d_0^3d_1d_2+28800d_0^3d_1+25920d_0^3d_2\\
   &+6720d_0^3d_3+480d_0^3d_4+2880d_0^2d_1^3+63360d_0^2d_1^2\\
   &+46080d_0^2d_1d_2+4320d_0^2d_1d_3+17280d_0^2d_1+2880d_0^2d_2^2\\
   &+46080d_0^2d_2+29280d_0^2d_3+6000d_0^2d_4+360d_0^2d_5\\
   &+23040d_0d_1^3+8640d_0d_1^2d_2+74880d_0d_1^2+123840d_0d_1d_2\\
   &+30720d_0d_1d_3+2160d_0d_1d_4+21600d_0d_2^2+3840d_0d_2d_3\\
   &+17280d_0d_2+30720d_0d_3+14640d_0d_4+2400d_0d_5+120d_0d_6\\
   &+720d_1^4+30720d_1^3+28800d_1^2d_2+2160d_1^2d_3+14400d_1^2\\
   &+3600d_1d_2^2+66240d_1d_2+36960d_1d_3+6720d_1d_4+360d_1d_5\\
   &+25200d_2^2+12000d_2d_3+840d_2d_4+480d_3^2+5760d_3\\
   &+7392d_4+2720d_5+360d_6+15d_7\bigr).
\end{aligned}
\end{equation}

Its top term is $15d_7/5760=d_7/384$, and $(-1)^4/(2^4\cdot4!)=1/384$: a
fourth independent check of \cref{thm:leading}.

\subsection{The same coefficients in the derivatives of \texorpdfstring{$\Psi$}{Psi}}

Write $P_k=\Psi^{(k)}(u)$.  Then
\[
  A_1(u)=\frac1{24}+\frac1{2L}-\frac{P_2+P_1^2}{2L^2},
\]
\begin{equation}\label{eq:A2psi}
\begin{aligned}
  A_2(u)={}&\frac{P_1}{24L}
    +\frac{2P_1+1}{4L^2}
    -\frac{P_1^3+3P_1^2+3P_1P_2+3P_2+P_3}{6L^3}\\
   &+\frac{4P_1^2P_2+4P_1P_3+2P_2^2+P_4}{8L^4}.
\end{aligned}
\end{equation}
The corresponding form of $A_3$ is long.  Its \emph{linear part} is the
following; the omitted terms are quadratic and higher in the derivatives of
$\Psi$ and reach $8.9\times10^{-7}$ over a period, which is four decades
below the $6.7\times10^{-2}$ swing of $A_3$ but is not negligible if six
digits are wanted.
\begin{equation}\label{eq:A3psi}
\begin{aligned}
  A_3(u)={}&-\frac7{2880}-\frac1{24L}-\frac1{4L^2}+\frac1{6L^3}
    +\frac{P_1+P_2}{24L^2}
    +\frac{48P_1+24P_2-P_3}{48L^3}\\
   &-\frac{6P_2+9P_3+P_4}{12L^4}
    +\frac{3P_4+P_5}{12L^5}
    -\frac{P_6}{48L^6}
    +(\text{quadratic and higher in the }P_k).
\end{aligned}
\end{equation}

\subsection{The \texorpdfstring{$\Psi$}{Psi}-free parts}

\begin{center}
\begin{tabular}{@{}clr@{}}
\toprule
$j$ & $\Psi$-free part of $A_j$ & value\\
\midrule
$1$ & $\dfrac1{24}+\dfrac1{2L}$ & $0.7630141871$\\[0.6em]
$2$ & $\dfrac1{4L^2}$ & $0.5203422453$\\[0.6em]
$3$ & $-\dfrac7{2880}-\dfrac1{24L}-\dfrac1{4L^2}+\dfrac1{6L^3}$ & $-0.0824216430$\\
$4$ & $\dfrac{-L^3-72L^2-816L+144}{1152L^4}$ & $-1.7167954639$\\
\bottomrule
\end{tabular}
\end{center}

The order-two value is remarkably clean: the $\Psi$-free part of $A_2$ is
exactly $1/(4\log^2 2)$, with no rational term and no $1/L$ term at all.

\subsection{The resulting expansions}

\[
\begin{aligned}
  \log F(x)&=\mathcal M(x)+\bigO(1/\lambda),\\
  \log F(x)&=\mathcal M(x)+\frac{A_1(\lambda)}{\lambda}+\bigO(1/\lambda^2),\\
  \log F(x)&=\mathcal M(x)+\frac{A_1(\lambda)}{\lambda}
    +\frac{A_2(\lambda)}{\lambda^2}+\bigO(1/\lambda^3),\\
  \log F(x)&=\mathcal M(x)+\sum_{j=1}^{N-1}\frac{A_j(\lambda)}{\lambda^j}
    +\bigO(1/\lambda^N)\qquad\text{for every }N.
\end{aligned}
\]

Written out at order three, with $d_k=d_k(\lambda)$ and $\lambda=\lambda(x)$,
\[
\begin{aligned}
  \log F(x)={}&-\frac L2\lambda^2+\Bigl(1+\frac L2\Bigr)\lambda
    -\frac12\log\lambda+C_F+\Psi(\lambda)\\
   &-\frac{6d_0^2+12d_0+6d_1+1}{12\lambda}\\
   &+\frac{8d_0^3+12d_0^2d_1+12d_0^2+48d_0d_1+12d_0d_2
      +6d_1^2+24d_1+20d_2+3d_3}{24\lambda^2}
    +\bigO(\lambda^{-3}).
\end{aligned}
\]
The order-three statement is
\texttt{log\_fabius\_sub\_twoSaddleCorrections\_isBigO} in
\nolinkurl{FabiusSecondSaddleExpansion.lean}.  Before this note the
formalization stopped at order two, because $A_2$ was not known.

\subsection{The right endpoint}\label{sec:right}

Nothing separate is needed at $x\to1^-$.  The symmetry $F(x)+F(1-x)=1$,
which holds on all of $\R$ for the bounded Fabius function, gives
$1-F(1-x)=F(x)$ identically, so the deficiency at the right endpoint is the
value at the left endpoint and every expansion above transfers verbatim:
\[
  \log\bigl(1-F(1-x)\bigr)=\mathcal M(x)
    +\sum_{j=1}^{N-1}\frac{A_j(\lambda(x))}{\lambda(x)^j}
    +\bigO(\lambda(x)^{-N}),\qquad x\to0^+,
\]
with the same $\mathcal M$, the same $\lambda$ and the same coefficients.
In particular $1-F(y)$ is flat to infinite order at $y=1$, with exactly the
same periodic modulation.  The reason this is worth a sentence rather than a
section is that the symmetry is exact, not asymptotic; there is no error term
to propagate.

\section{The leading-derivative law}\label{sec:leading}

\begin{theorem}[Leading term]\label{thm:leading}
The coefficient of $d_{2j-1}$ in both $a_j$ and $A_j$ equals
$(-1)^j/(2^jj!)$, so the highest-derivative term of $A_j$ is
\[
  \frac{(-1)^j\,\Psi^{(2j)}(u)}{2^j\,j!\,L^{2j}}.
\]
\end{theorem}

\begin{proof}[Proof sketch]
$A_j$ reaches $\Psi^{(2j)}$ only through the jet $d_{2j-1}$, and $d_{2j-1}$
occurs in the expansion only linearly, through $E_{2j}$.  In $E_{2j}$ the jet
term is $\ii^{\,2j}d_{2j-1}v^{2j}/(2j)!=(-1)^jd_{2j-1}v^{2j}/(2j)!$; contracting
$v^{2j}$ against the Gaussian gives $(2j-1)!!$; and $(2j)!=2^jj!\,(2j-1)!!$.
Nothing else at order $2j$ is linear in $d_{2j-1}$, and the formal logarithm
does not change a linear term.  Finally the top term of $d_{2j-1}$ is
$\Psi^{(2j)}/L^{2j}$ because $c(n,n)=1$.
\end{proof}

Checks: $j=1$ gives $-\Psi''/(2L^2)$, $j=2$ gives $+\Psi''''/(8L^4)$ and $j=3$
gives $-\Psi^{(6)}/(48L^6)$ --- exactly the top terms displayed in
\cref{eq:A2psi,eq:A3psi}.

\section{Fourier form and the amplitude law}\label{sec:fourier}

Because $\Psi^{(m)}(u)=\sum_{k\ne0}(2\pi\ii k)^m\Psihat(k)\e^{2\pi\ii ku}$,
\cref{sec:general} together with \cref{eq:psihat} turns every coefficient into
an explicit Gamma--zeta Fourier object.  Discarding the terms of degree two
and higher in the derivatives of $\Psi$ leaves the linearized Fourier form
\[
  A_j(u)=\langle A_j\rangle
   +2\operatorname{Re}\Bigl(\sum_{k\ge1}c_j(k)\Psihat(k)\e^{2\pi\ii ku}\Bigr)
   +R_j(u),
\]
where the discarded remainder is measured at
$\max_u|R_j(u)|=1.6\times10^{-8}$ for $j=2$ and $8.9\times10^{-7}$ for
$j=3$, growing with $j$ like everything else here.  It is negligible against
the oscillation it sits inside, and it is not negligible in absolute terms;
\cref{sec:general} records the same caution for the means.
with $c_j(k)$ the polynomial in $w=2\pi\ii k$ read off from
\cref{eq:A2psi,eq:A3psi}.  For the first two,
\[
  c_1(k)=-\frac{w^2}{2L^2},
  \qquad
  c_2(k)=\frac{w}{24L}+\frac{w}{2L^2}-\frac{3w^2+w^3}{6L^3}+\frac{w^4}{8L^4}.
\]
Consequences for $A_2$, all confirmed numerically in \cref{sec:numerics}.
The left column is what the linearized form gives, the right what the exact
$A_2$ gives; they are quoted separately because the difference is real, if
small, and quoting ten digits of a linearization would be an overclaim.
\[
\begin{aligned}
  \text{mean}&&\tfrac1{4L^2}=0.5203422452514&\qquad
    \langle A_2\rangle=0.5203422413049\\
  k=1\ \text{amplitude}&&2|c_2(1)\Psihat(1)|=1.9398782\times10^{-3}&\qquad
    \text{peak-to-peak }3.8797564\times10^{-3}\\
  \text{maximum at}&&\{\lambda\}=0.1011297&\qquad\{\lambda\}=0.1011281\\
  \text{minimum at}&&\{\lambda\}=0.6011297&\qquad\{\lambda\}=0.6011313
\end{aligned}
\]
The exact peak-to-peak is $3.8797564\times10^{-3}$, agreeing with the
linearized amplitude to eight digits, and the extrema agree to six.  The
maximum of the exact $A_2$ is $0.522282117$ and its minimum is
$0.518402361$.

\subsection{The amplitude law}

The top term of \cref{thm:leading} dominates the oscillation, and it multiplies
the $k$-th Fourier mode by $(2\pi k)^{2j}$.  Taking $k=1$,
\begin{equation}\label{eq:amplitude}
  \text{amplitude of }A_j\ \approx\
  2|\Psihat(1)|\,\frac{(2\pi^2/L^2)^j}{j!},
  \qquad \frac{2\pi^2}{L^2}=41.0845769104 .
\end{equation}

\begin{center}
\begin{tabular}{@{}cr@{}}
\toprule
$j$ & predicted $k=1$ amplitude\\
\midrule
$1$ & $8.73\times10^{-5}$\\
$2$ & $1.79\times10^{-3}$\\
$3$ & $2.46\times10^{-2}$\\
$4$ & $2.52\times10^{-1}$\\
$5$ & $2.07$\\
$6$ & $1.42\times10^{1}$\\
$7$ & $8.33\times10^{1}$\\
\bottomrule
\end{tabular}
\end{center}

These are estimates from the top term alone.  Comparing them with the truth
requires care about a convention: an \emph{amplitude} is half a
peak-to-peak swing, and mixing the two silently doubles every second
entry.  Everything below is an amplitude.

\begin{center}
\begin{tabular}{@{}crrr@{}}
\toprule
$j$ & \cref{eq:amplitude} & actual & ratio\\
\midrule
$1$ & $8.73\times10^{-5}$ & $8.73\times10^{-5}$ & $1.00$\\
$2$ & $1.79\times10^{-3}$ & $1.94\times10^{-3}$ & $1.08$\\
$3$ & $2.46\times10^{-2}$ & $\approx3\times10^{-2}$ & $\approx1.2$\\
$4$ & $2.52\times10^{-1}$ & $\approx0.41$ & $\approx1.6$\\
\bottomrule
\end{tabular}
\end{center}

At $j=1$ the top term is the only term, so the estimate is exact.  At
$j=2$ the $w^3$ term of $c_2$ adds in phase and lifts the amplitude above
the estimate; the entries at $j=3$ and $j=4$ are halves of the swings
measured in \cref{sec:numerics}.  The ratio drifts upward because the
subleading terms of $c_j$ contribute more as $j$ grows.  So
\cref{eq:amplitude} gives the growth rate and not the constant, and should
never be used to size a particular coefficient.  So the periodic parts of
the coefficients grow rapidly with $j$, even though $\Psi$ itself oscillates by
only $2\times10^{-6}$: differentiation progressively cancels the smallness of
$\Psihat(1)$.  Any numerical extraction of a coefficient beyond the first must
therefore control the phase; see \cref{sec:traps}.

\subsection{Why the expansion is only asymptotic}

$|\Psihat(k)|$ decays like $\exp(-\pi|\chi_k|/2)=\exp(-\pi^2k/L)$, that is,
about $\e^{-14.24k}$, while mode $k$ is amplified by
$(2\pi k)^{2j}/(2^jj!\,L^{2j})$.  The $k$-th contribution therefore behaves like
$\exp(2j\log k-14.24k)$, maximized near $k^\ast=2j/14.24=0.14j$.  For $j$ beyond
about $8$ the higher modes dominate, and $\sup|A_j|$ grows faster than
geometrically.

What is worth noticing is \emph{which} divergence this is.  A saddle-point
expansion normally diverges factorially, because the Gaussian moments grow
like $(2j-1)!!$ against a fixed exponent.  Here that growth is exactly
cancelled: the $(2j)!$ in $E_{2j}$ absorbs it, which is the content of
\cref{thm:leading} and the reason the coefficient of the top jet is the mild
$1/(2^jj!)$.  The divergence that remains is of a different origin
altogether --- it is driven by the high-frequency content of $\Psi$, through
the $(2\pi k)^{2j}$ that differentiation applies to the $k$-th mode.  A
reader who expects the usual factorial mechanism will look for it in the
Gaussian integration, where it is not.  (This framing is due to the parallel
write-up of the main article, where it also appears.)

\begin{remark}
The expansion is a Poincar\'e asymptotic expansion: each truncation is valid
with an $\bigO(\lambda^{-N})$ remainder for fixed $N$, and there is no claim,
here or in the formalization, that the series converges for any fixed $x$.
The estimate in this subsection is a heuristic size estimate, not a theorem.
\end{remark}

\section{Numerical verification}\label{sec:numerics}

$F(2^{-n})$ is computed \emph{exactly} as a rational number from the
inverse-power recursion used by the Lean evaluator
(\nolinkurl{Arithmetic.lean}, \lstinline[style=pseudo]{fabiusInversePowTwoTable}):
\[
  v_0=1,\qquad
  v_{n+1}=\frac{1}{2^{\,n+1}-1}
    \sum_{k=0}^{n}\frac{v_k}{2^{\binom{n+1}{2}-\binom{k}{2}}\,(n-k+2)!},
  \qquad F(2^{-n})=v_n.
\]
Sanity checks: $v_1=1/2=F(1/2)$ and $v_2=5/72=F(1/4)$.  All computations below
use $260$-digit working precision, with $\Psi$ and its derivatives summed over
six Gamma--zeta modes and $\zeta$ evaluated by Euler--Maclaurin
(\cref{sec:traps}).

Define the successive residuals
\[
  r_1=\lambda\bigl(\log F-\mathcal M\bigr),\quad
  r_2=\lambda^2\Bigl(\log F-\mathcal M-\frac{A_1}{\lambda}\Bigr),\quad
  r_3=\lambda^3\Bigl(\log F-\mathcal M-\frac{A_1}{\lambda}-\frac{A_2}{\lambda^2}\Bigr).
\]

\begin{center}\small
\begin{tabular}{@{}rrrrrrrr@{}}
\toprule
$n$ & $\lambda$ & $r_1$ & $A_1(\lambda)$ & $r_2$ & $A_2(\lambda)$ & $r_3$ & $A_3(\lambda)$\\
\midrule
$10$ & $13.78503056$ & $0.7997155875$ & $0.7629678468$ & $0.5065687288$ & $0.5195595404$ & $-0.1790787344$ & $-0.0868938679$\\
$20$ & $24.62186833$ & $0.7836542295$ & $0.7629268731$ & $0.5103462408$ & $0.5184188039$ & $-0.1987615843$ & $-0.1111875606$\\
$40$ & $45.50804986$ & $0.7742689337$ & $0.7629490545$ & $0.5151456258$ & $0.5187247759$ & $-0.1628801414$ & $-0.1121601297$\\
$60$ & $66.04538587$ & $0.7709808331$ & $0.7630910552$ & $0.5210834261$ & $0.5221643489$ & $-0.0713899656$ & $-0.0512594569$\\
$80$ & $86.43351899$ & $0.7689739453$ & $0.7629823191$ & $0.5178773369$ & $0.5193822939$ & $-0.1300787245$ & $-0.1046130809$\\
$120$ & $126.9885547$ & $0.7671769004$ & $0.7630717251$ & $0.5213102874$ & $0.5218167515$ & $-0.0643151389$ & $-0.0540815552$\\
$160$ & $167.3870441$ & $0.7661091776$ & $0.7630070725$ & $0.5192522062$ & $0.5199082072$ & $-0.1098060709$ & $-0.0973365576$\\
\bottomrule
\end{tabular}
\end{center}

The decisive test is that each residual, minus its predicted coefficient, is
$\bigO(1/\lambda)$.  Writing $e_j=r_j-A_j$:

\begin{center}\small
\begin{tabular}{@{}rrrrr@{}}
\toprule
$n$ & $\{\lambda\}$ & $\lambda e_1$ & $\lambda e_2$ & $\lambda e_3$\\
\midrule
$100$ & $0.737929$ & $0.51802529$ & $-0.11213991$ & $-1.7523898$\\
$112$ & $0.893526$ & $0.52019699$ & $-0.07798492$ & $-1.3773605$\\
$124$ & $0.033795$ & $0.52164121$ & $-0.06156984$ & $-1.3168268$\\
$136$ & $0.161500$ & $0.52166737$ & $-0.06826635$ & $-1.5329675$\\
$148$ & $0.278716$ & $0.52062406$ & $-0.08861040$ & $-1.8452719$\\
$160$ & $0.387044$ & $0.51925221$ & $-0.10980607$ & $-2.0872350$\\
\bottomrule
\end{tabular}
\end{center}

$\lambda e_1$ reproduces $A_2$ with mean $0.5203$; $\lambda e_2$ reproduces
$A_3$ with mean $-0.0824$; and $\lambda e_3$ is bounded, locating $A_4$ between
about $-1.3$ and $-2.1$.  Each column stays bounded and oscillates with the
phase rather than drifting, which is precisely the assertion that the
corresponding coefficient is correct and the next one is $\bigO(1)$.

Pointwise, $A_2$ is confirmed at the $10^{-5}$ level once the next order is
included.  For $140\le n\le160$:

\begin{center}\small
\begin{tabular}{@{}rrrrr@{}}
\toprule
$n$ & $\{\lambda\}$ & $A_2(\lambda)$ & $r_2$ & $r_2-A_3(\lambda)/\lambda$\\
\midrule
$140$ & $0.201650$ & $0.5219078895$ & $0.5214039042$ & $0.5218323910$\\
$145$ & $0.250301$ & $0.5214906179$ & $0.5209456288$ & $0.5214143255$\\
$150$ & $0.297351$ & $0.5209853273$ & $0.5203979284$ & $0.5209087825$\\
$155$ & $0.342900$ & $0.5204425007$ & $0.5198167900$ & $0.5203665207$\\
$160$ & $0.387044$ & $0.5199082072$ & $0.5192522062$ & $0.5198337121$\\
\bottomrule
\end{tabular}
\end{center}

The last two columns agree to about $7.5\times10^{-5}$, which is
$A_4/\lambda^2$ with $A_4\approx-2$.  Note that the \emph{variation} of $A_2$
over this window is $2.0\times10^{-3}$, twenty-six times the residual, so this
is a genuine pointwise verification of the periodic shape and not merely of a
mean.

\begin{remark}
An independent high-precision computation performed in a separate worktree,
using half-moments to $n=400$ at $400$-digit precision and a different route to
$\lambda$, located the minimum of the order-$\lambda^{-2}$ residual at
$\{\lambda\}=0.602$ and measured a peak-to-peak spread of about
$3\times10^{-3}$; \cref{eq:A2psi} predicts a minimum at $0.6011313$ and a
peak-to-peak of $3.8798\times10^{-3}$ over the full period.  That is a phase
prediction, not a magnitude fit, and it was made before the comparison.
\end{remark}

\subsection{The fourth coefficient}

The same apparatus confirms \cref{eq:A4jets}.  Writing
$r_4=\lambda^4(\log F-\mathcal M-A_1/\lambda-A_2/\lambda^2-A_3/\lambda^3)$,
at $300$-digit precision with eight Gamma--zeta modes:

\begin{center}\small
\begin{tabular}{@{}rrrr@{}}
\toprule
$n$ & $\{\lambda\}$ & $r_4$ & $A_4(\lambda)$\\
\midrule
$120$ & $0.988555$ & $-1.299548008$ & $-1.305123624$\\
$136$ & $0.161500$ & $-1.532967481$ & $-1.506670619$\\
$152$ & $0.315745$ & $-1.939809221$ & $-1.890185255$\\
$160$ & $0.387044$ & $-2.087234974$ & $-2.035031382$\\
$176$ & $0.519792$ & $-2.166743696$ & $-2.124617975$\\
$192$ & $0.641266$ & $-1.992935176$ & $-1.970225155$\\
$200$ & $0.698346$ & $-1.853347748$ & $-1.840001848$\\
\bottomrule
\end{tabular}
\end{center}

$\lambda(r_4-A_4)$ stays inside $[-8.8,0.8]$ over $120\le n\le200$ with no
drift, which places $A_5$ and confirms $A_4$.  The $\Psi$-free part
$-1.7167954639$ sits inside the measured range, and the swing of the
$A_4$ column, $0.82$ peak to peak, is $1.6$ times the leading-term estimate
$0.50$ of \cref{eq:amplitude} --- the same upward drift of that ratio seen
at $j=2$ ($1.08$) and $j=3$ ($1.24$), because the subleading terms of
$c_j(1)$ contribute more as $j$ grows.  \Cref{eq:amplitude} is a
leading-term estimate and should be read as one.

\subsection{A worked example}\label{sec:worked}

It is worth seeing the expansion used once, on a value that can be written
down exactly.  At $x=2^{-10}$ the evaluator gives
\[
  F(2^{-10})=\frac{134\,926\,369}{1\,246\,394\,851\,358\,539\,387\,238\,350\,848\,000},
  \qquad \log F(2^{-10})=-50.57756828234216081254\dots,
\]
and $\lambda=13.7850305606$, so $\{\lambda\}=0.785031$.  The successive
truncations miss $\log F$ by

\begin{center}\small
\begin{tabular}{@{}lrr@{}}
\toprule
truncation & error at $n=10$ & error at $n=160$\\
\midrule
$\mathcal M(x)$ & $5.80\times10^{-2}$ & $4.58\times10^{-3}$\\
${}+A_1/\lambda$ & $2.67\times10^{-3}$ & $1.85\times10^{-5}$\\
${}+A_2/\lambda^2$ & $-6.84\times10^{-5}$ & $-2.34\times10^{-8}$\\
${}+A_3/\lambda^3$ & $-3.52\times10^{-5}$ & $-2.66\times10^{-9}$\\
\bottomrule
\end{tabular}
\end{center}

where $n=160$ has $\lambda=167.387044$ and
$\log F(2^{-160})=-9489.629874649676879065\dots$

Two things are visible.  At $n=160$ each order gains roughly a factor
$\lambda$, as it should: $250$, then $790$.  The last gain is only $8.8$,
and that is not a defect --- the third-order error is already dominated by
$A_4/\lambda^4$, which \cref{eq:A4jets} puts at
$-1.72/167.387^4=-2.19\times10^{-9}$ against a measured
$-2.66\times10^{-9}$.  So the table is measuring $A_4$, not a failure.

At $n=10$, by contrast, the fourth term barely helps: the error falls only
from $6.8\times10^{-5}$ to $3.5\times10^{-5}$.  With $\lambda=13.8$ and
$A_4\approx-1.7$, the term $A_4/\lambda^4$ is itself $4.7\times10^{-5}$,
comparable to the error it is meant to reduce.  This is the ordinary
behaviour of a Poincar\'e expansion near the end of its useful range, and it
is the practical reason \cref{sec:fourier} declines to claim convergence:
the coefficients grow, so for a fixed $x$ there is a best truncation, and
going past it makes the approximation worse.  At $x=2^{-10}$ that optimum is
already at hand around the third term.

\subsection{An independent confirmation to $n=400$}

The following was produced in a separate worktree, by a different agent,
from a different representation of $F$, and is reproduced with its
permission.  Half-moments were computed to $n=400$ at $400$-digit working
precision from
\[
  d_n=\frac1{(n+1)(2^n-1)}\sum_{k<n}\binom{n+1}{k}d_k,
\]
whose terms are all positive, so there is no cancellation; the values agree
with the exact rationals to $400$ digits at $n=5,20,60$.  Then
$\log F(2^{-n})=\log d_n-\log n!-\binom n2\log2$, with $\lambda$ from the
lower Lambert branch (satisfying $\lambda2^{-\lambda}=x$ to $10^{-21}$) and
$\Psi$ from ten Gamma--zeta modes.  Nothing in this table shares an
implementation with \cref{sec:numerics}.

\begin{center}\small
\begin{tabular}{@{}rrrrrr@{}}
\toprule
$n$ & $\lambda$ & $\{\lambda\}$ & $r_2$ measured & $A_2(\lambda)$ & $r_3$\\
\midrule
$100$ & $106.73793$ & $0.7379$ & $0.51802529$ & $0.51907590$ & $-0.11214$\\
$125$ & $132.04488$ & $0.0449$ & $0.52169792$ & $0.52216224$ & $-0.06131$\\
$150$ & $157.29735$ & $0.2974$ & $0.52039793$ & $0.52098533$ & $-0.09240$\\
$175$ & $182.51185$ & $0.5119$ & $0.51801879$ & $0.51869969$ & $-0.12427$\\
$200$ & $207.69835$ & $0.6984$ & $0.51821775$ & $0.51875313$ & $-0.11120$\\
$225$ & $232.86334$ & $0.8633$ & $0.52015631$ & $0.52049092$ & $-0.07792$\\
$250$ & $258.01129$ & $0.0113$ & $0.52175813$ & $0.52198119$ & $-0.05755$\\
$275$ & $283.14540$ & $0.1454$ & $0.52199208$ & $0.52220755$ & $-0.06101$\\
$300$ & $308.26804$ & $0.2680$ & $0.52104845$ & $0.52130958$ & $-0.08050$\\
$325$ & $333.38103$ & $0.3810$ & $0.51967269$ & $0.51997997$ & $-0.10244$\\
$350$ & $358.48577$ & $0.4858$ & $0.51856496$ & $0.51889001$ & $-0.11653$\\
$375$ & $383.58340$ & $0.5834$ & $0.51810490$ & $0.51841439$ & $-0.11872$\\
$400$ & $408.67481$ & $0.6748$ & $0.51833670$ & $0.51860654$ & $-0.11028$\\
\bottomrule
\end{tabular}
\end{center}

The decisive column is the last.  Over this range $\lambda$ quadruples, so
an error of even $10^{-4}$ in $A_2$ would appear as a drift in $r_3$ of
order $3\times10^{-2}$ with a definite sign.  Instead $r_3$ is $-0.112$ at
$n=100$ and $-0.110$ at $n=400$, never leaving $[-0.125,-0.058]$ in
between, and oscillating rather than drifting --- which is what a periodic
$A_3$ requires.  Its sample mean over the thirteen rows is about $-0.088$,
against the $\Psi$-free part $-0.0824216430$ of $A_3$; the sample is not a
uniform phase average and carries an $A_4/\lambda$ term, so the gap is
expected.

The $A_2(\lambda)$ column also straddles the predicted mean
$1/(4L^2)=0.5203422452514019$, and the row-by-row differences
$r_2-A_2$ match the prediction $A_3(\lambda)/\lambda$ to about $10^{-6}$:
at $\{\lambda\}=0.2680$ the predicted difference is $-2.62\times10^{-4}$
against a measured $-2.611\times10^{-4}$, and at $\{\lambda\}=0.6748$,
$-2.71\times10^{-4}$ against $-2.698\times10^{-4}$.

Neither this table nor \cref{sec:numerics} is machine-checked.  What is
machine-checked is listed in \cref{sec:lean}.

\section{Map to the Lean development}\label{sec:lean}

\begin{center}\small
\begin{longtable}{@{}lll@{}}
\toprule
object & Lean name & module\\
\midrule
\endfirsthead
\toprule
object & Lean name & module\\
\midrule
\endhead
$\lambda(x)$ & \texttt{fabiusLambertPhase} & \texttt{FabiusLambertPhase.lean}\\
$P$ & \texttt{negativeLaplacePeriodicCorrection} & \texttt{PeriodicCorrection.lean}\\
$\Psi$ & \texttt{negativeLaplacePsi} & \texttt{PeriodicCorrection.lean}\\
$\Psihat(k)$ & Fourier-coefficient theorems & \texttt{PeriodicFourier.lean}\\
$C_F$ & \texttt{fabiusSharpAsymptoticConstant} & \texttt{FabiusSharpConstant.lean}\\
$\mathcal M(x)$ & \texttt{fabiusSharpLambertMain} & \texttt{FabiusSharpLambertMain.lean}\\
$s_n$ & \texttt{negativeLaplaceJetSlope} & \texttt{FabiusSaddleCoefficientRecurrence.lean}\\
$p_n$ & \texttt{negativeLaplacePeriodicJet} & \texttt{FabiusSaddleCoefficientRecurrence.lean}\\
$d_n$ & \texttt{negativeLaplaceBoundedExponentJet} & \texttt{FabiusSaddleCoefficientRecurrence.lean}\\
$c(n,m)$ & \texttt{negativeLaplaceJetStirling} & \textbf{\texttt{FabiusSaddleJetClosedForm.lean}}\\
$\Psi^{(m+1)}/L^{m+1}$ & \texttt{negativeLaplacePsiScaledDeriv} & \textbf{\texttt{FabiusSaddleJetClosedForm.lean}}\\
closed-form $p_n$ & \texttt{negativeLaplacePeriodicJet\_eq\_closedForm} & \textbf{\texttt{FabiusSaddleJetClosedForm.lean}}\\
closed-form $d_n$ & \texttt{negativeLaplaceBoundedExponentJet\_eq\_closedForm} & \textbf{\texttt{FabiusSaddleJetClosedForm.lean}}\\
$E_m$ & \texttt{negativeLaplaceExponentCoefficient} & \texttt{FabiusSaddleCoefficientRecurrence.lean}\\
closed-form $E_m$ & \texttt{negativeLaplaceExponentCoefficient\_succ\_eq} & \textbf{\texttt{FabiusSaddleExponentClosedForm.lean}}\\
$\operatorname{expCoeff}$ & \texttt{SaddleExpansion.expCoeff} & \texttt{SaddleExpansionAlgebra.lean}\\
contraction & \texttt{SaddleExpansion.gaussianPolynomialContraction} & \texttt{GaussianPolynomialContraction.lean}\\
$a_j$ & \texttt{fabiusSaddleMassCoefficient} & \texttt{FabiusSaddleExpansionCoefficients.lean}\\
$\operatorname{logCoeff}$ & \texttt{SaddleExpansion.logCoeff} & \texttt{SaddleLogExpansionAlgebra.lean}\\
$A_j$ & \texttt{fabiusSaddleLogCoefficient} & \texttt{FabiusSaddleExpansionCoefficients.lean}\\
$A_1$ & \texttt{fabiusFirstSaddleCorrection} & \texttt{FabiusFirstSaddleCorrection.lean}\\
$A_2$ & \texttt{fabiusSecondSaddleCorrection} & \textbf{\texttt{FabiusSecondSaddleCorrection.lean}}\\
$a_2$ & \texttt{fabiusSaddleMassCoefficient\_two} & \textbf{\texttt{FabiusSecondSaddleCorrection.lean}}\\
$c(n,m)$ identified & \texttt{negativeLaplaceJetStirling\_eq\_coeff} & \textbf{\texttt{FabiusSaddleJetStirling.lean}}\\
$\operatorname{expCoeff}$ at $3,4$ & \texttt{expCoeff\_three}, \texttt{expCoeff\_four} & \textbf{\texttt{FabiusSecondSaddleCorrection.lean}}\\
order-three expansion & \texttt{log\_fabius\_sub\_twoSaddleCorrections\_isBigO} & \textbf{\texttt{FabiusSecondSaddleExpansion.lean}}\\
the expansion & \texttt{log\_fabius\_sub\_sharpLambertMain\_hasAsymptoticExpansion} & \texttt{FabiusFullAsymptoticExpansion.lean}\\
truncation & \texttt{log\_fabius\_sub\_sharpLambertExpansion\_isBigO} & \texttt{FabiusFullAsymptoticExpansion.lean}\\
exact $F(2^{-n})$ & \texttt{fabiusInversePowTwoTable}, \texttt{evalFabiusDyadic} & \texttt{Arithmetic.lean}\\
\bottomrule
\end{longtable}
\end{center}

Modules in bold are the ones added with this document.

\subsection{What is machine-checked, and what is not}

The three kinds of evidence in this note are not interchangeable, so they
are separated here.

\emph{Machine-checked.}  The following modules compile under Lean 4.32.0
with Mathlib \texttt{v4.32.0}, with no \texttt{sorry}, no declared axiom and
no \texttt{opaque}, on top of a serialized topological build of their whole
import closure:
\begin{itemize}
\item \texttt{FabiusSaddleJetClosedForm.lean}
\item \texttt{FabiusSaddleExponentClosedForm.lean}
\item \texttt{FabiusSaddleJetStirling.lean}
\item \texttt{FabiusSaddleLeadingCoefficient.lean}
\item \texttt{FabiusSecondSaddleCorrection.lean}
\end{itemize}

\emph{Written but not yet compiled} at this snapshot:
\begin{itemize}
\item \texttt{FabiusSecondSaddleExpansion.lean}
\end{itemize}
That one module substitutes $A_2$ into the truncated expansion and states
$\log F(x)=\mathcal M(x)+A_1/\lambda+A_2/\lambda^2+\bigO(\lambda^{-3})$ for
every bounded Fabius function.  Its import closure is the full asymptotic
aggregate, about $110$ modules, which was beyond the build budget of this
snapshot; a reader should treat it as a claim until the build log says
otherwise.  Everything it depends on is machine-checked, and its own proofs
are three rewrites.

\emph{Verified symbolically but not formalized.}  \Cref{thm:mass,thm:log},
the composition formulas, were checked against the original recursions in
\textsc{sympy} for $j\le3$, with residuals simplifying to identically zero.
The coefficients $A_3$ and $A_4$ of \cref{eq:A3jets,eq:A4jets} were derived
by the same pipeline, which reproduces $A_1$ exactly against the
independently formalized \texttt{fabiusFirstSaddleCorrection}.  That
agreement is the pipeline's calibration, not a proof of $A_3$ or $A_4$.

\emph{Verified numerically only.}  Everything in \cref{sec:numerics}.
High-precision agreement between an exact rational evaluation of $F$ and a
closed form is strong evidence and is not a proof; in particular it cannot
distinguish a correct coefficient from one differing by a quantity below the
noise floor of the sampled range.

No claim in this note rests on the numerics alone except the location of
$A_5$, which is stated only as a bound on an observed residual.


\section{What is not claimed}\label{sec:notclaimed}

\begin{itemize}
\item \textbf{No elementary all-orders form.}  The coefficients are evaluated at
  the exact phase $\lambda(x)$, not at a truncation of it.  The phase itself has
  a separate all-orders expansion in $S=-\log x$ and $z=\log S-\log L$,
  \[
    L\lambda=S+z+\sum_{j<N}\frac{p_j(z)}{S^j}+\bigO\!\left(\frac{(1+|z|)^N}{S^N}\right),
    \quad p_1=z,\quad p_2=z-\frac{z^2}{2},\quad p_3=z-\frac{3z^2}{2}+\frac{z^3}{3},
  \]
  with the closed Stirling form
  $p_j(z)=\sum_{m=1}^{j}\frac{(-1)^{m+1}}{m!}s(j,j-m+1)z^m$.  But
  $\mathcal M(x)$ contains $-\tfrac L2\lambda^2$, so an error $\delta$ in
  $\lambda$ produces an error of order $L\lambda\delta$ in $\log F$; and the
  periodic coefficients are evaluated at $\lambda$, where a truncation error
  moves the phase.  Composing the two expansions therefore needs its own
  theorem, which is deliberately not asserted.
\item \textbf{No convergence}; see \cref{sec:fourier}.
\item \textbf{No effective constants.}  The $\bigO(\lambda^{-N})$ remainders are
  proved, but their implied constants are not surfaced.
\item \textbf{No claim that any $A_j$ is nowhere zero.}  $A_1$, $A_2$ and $A_3$
  are identified; the theorems identify the coefficient functions, they do not
  assert that these functions have no zeros.
\end{itemize}

\section{Two traps}\label{sec:traps}

\begin{warning}[The phase trap]
Sampling $n=10,20,40,80,160,320$ to estimate a coefficient is invalid.  Doubling
$n$ adds about $1$ to $\log_2\lambda$, so $\{\lambda\}$ returns to nearly the
same place and a geometric sample holds the phase almost fixed.  Since every
coefficient is one-periodic, such a sample cannot see the oscillation at all;
it shows only the $\bigO(1/\lambda)$ drift in magnitude, and looks like
convergence to a constant.  This mistake was made independently on this project
before it was caught.  Always sweep $n$ densely and report against
$\{\lambda\}$, as in \cref{sec:numerics}.
\end{warning}

\begin{warning}[The zeta trap]
Many $\zeta$ implementations route through the alternating Dirichlet eta
series, dividing by $1-2^{1-s}$.  At $s=1-\chi_k$ one has
$2^{1-s}=\e^{2\pi\ii k}=1$ exactly, so that prefactor has a pole at precisely
the arguments \cref{eq:psihat} needs.  The singularity is removable in $\zeta$
but not in that formula, and the result is either an overflow or a plausible
wrong number.  Concretely, \texttt{mpmath.zeta} at these arguments returns
values whose seventh significant digit drifts with the working precision:
\[
\begin{aligned}
  \texttt{mp.dps}=30:&\quad \Psihat(1)=7.55865005\times10^{-7}-7.47011001\times10^{-7}\,\ii,\\
  \texttt{mp.dps}=50:&\quad \Psihat(1)=7.55864674\times10^{-7}-7.47009636\times10^{-7}\,\ii,\\
  \texttt{mp.dps}=80:&\quad \Psihat(1)=7.55865060\times10^{-7}-7.47008702\times10^{-7}\,\ii,
\end{aligned}
\]
whereas a hand-rolled Euler--Maclaurin $\zeta$ ($N=300$, $25$ Bernoulli terms)
is precision-stable and returns
$7.55864754097769\times10^{-7}-7.47010355898008\times10^{-7}\,\ii$ at every
working precision.  Use Euler--Maclaurin or Hurwitz, and always check at two
precisions.
\end{warning}

\section{Notation crosswalk}\label{sec:crosswalk}

Three namings for the same objects appear across this project.

\begin{center}
\begin{tabular}{@{}lll@{}}
\toprule
this note / the article & Wikipedia draft & Lean\\
\midrule
$\lambda(x)$ & $\lambda$ & \texttt{fabiusLambertPhase x}\\
$\Psi$ & $\Psi$ & \texttt{negativeLaplacePsi}\\
$\mathcal M(x)$ & main term & \texttt{fabiusSharpLambertMain x}\\
$A_j$ & $P_j$ & \texttt{fabiusSaddleLogCoefficient j}\\
$a_j$ & $b_j$ & \texttt{fabiusSaddleMassCoefficient j}\\
$A_1$ & $P_1$ & \texttt{fabiusFirstSaddleCorrection}\\
$A_2$ & $P_2$ & \texttt{fabiusSecondSaddleCorrection}\\
$C_F$ & $C_F$ & \texttt{fabiusSharpAsymptoticConstant}\\
\bottomrule
\end{tabular}
\end{center}

The Wikipedia draft calls the coefficients merely continuous; they are in fact
$C^\infty$, one-periodic and bounded, and \cref{sec:general} shows that they are
polynomials in smooth periodic functions.

The jets and the exponent coefficients have no counterpart in the Wikipedia
draft, and in the main article they carry different letters, because $p$, $d$
and $E$ are all taken there: $p_j$ is the Lambert phase polynomial and $p_m$
the binary prefix, $d_n$ is the half moment in some sixty places, and $E$ is
both the tail error and the exponential quotient.  The article's names are:

\begin{center}
\begin{tabular}{@{}lll@{}}
\toprule
this note & the main article & Lean\\
\midrule
$p_n$ & $Z_n$ & \texttt{negativeLaplacePeriodicJet}\\
$d_n$ & $\widetilde Z_n$ & \texttt{negativeLaplaceBoundedExponentJet}\\
$E_m$ & $\Xi_m$ & \texttt{negativeLaplaceExponentCoefficient}\\
$\dd/\dd t$ & $\mathrm D$ & ---\\
\bottomrule
\end{tabular}
\end{center}

The tilde in $\widetilde Z_n$ is apt: the two jet families differ by the
single constant $(-1)^{n+1}n!$.

\section{Open items}\label{sec:open}

\begin{enumerate}
\item Formalize \cref{thm:mass,thm:log} for the generic
  $\operatorname{expCoeff}$ and $\operatorname{logCoeff}$ recursions.  This is
  the one part of the general formula that remains unformalized, so it is
  worth recording what it would actually take.  Mathlib has
  \texttt{Composition n} with a \texttt{Fintype} instance, but no first-block
  decomposition equivalence
  $\mathcal C(n+1)\simeq\coprod_{b=1}^{n+1}\mathcal C(n+1-b)$, which is what
  an induction against the recursion needs.

  The route that avoids \texttt{Composition} altogether is to work with the
  $k$-fold convolution powers $E^{*k}$ directly, as plain finite sums, and
  prove $\operatorname{expCoeff}(E)_n=\sum_{k\le n}(k!)^{-1}(E^{*k})_n$ when
  $E_0=0$.  The pleasant part is that the product rule is a triviality at the
  level of coefficients: writing $(\mathrm DA)_n=nA_n$,
  $\mathrm D(A*B)=(\mathrm DA)*B+A*(\mathrm DB)$ is just $n=j+(n-j)$ inside
  the convolution sum.  The unpleasant part is that promoting it to
  $\mathrm D(E^{*(k+1)})=(k+1)\,E^{*k}*\mathrm DE$ needs associativity of
  convolution, which is a double-sum reindexing that Mathlib does not state
  for bare sequences either.  Either ingredient is a self-contained piece of
  \texttt{Finset} combinatorics; neither is deep, and both are larger than the
  theorem they support.
\item Formalize \cref{thm:leading} for general $j$.
\item Surface effective constants in the $\bigO(\lambda^{-N})$ remainders.
\item Prove or refute convergence of the series in \cref{sec:fourier}.
\item Compute $A_4$ in closed form; the numerics of \cref{sec:numerics} place
  its mean near $-1.7$ and \cref{eq:amplitude} predicts an oscillation
  amplitude of about $0.25$.  The machinery is in place: it needs
  $\operatorname{expCoeff}$ at order $6$, the jets $d_0,\dots,d_5$, and
  Gaussian moments to $v^{18}$.
\end{enumerate}

\end{document}
